\chapter{Assignment 5.1: Model Organism}

For what I understood also at the lecture, a model organism is a non-human organism that is studied to understand biological processes. 
The idea comes first from the situation that we cannot experiment on humans, 
so we started to seek organisms that are compatible with us, because, 
at the end of the story we all share a common evolutionary path, 
this means that the cells that are living in animals is often similar on how it works in us.

Also, besides the fact that we cannot test on humans, we use model organisms because they are rapid to grow (for example mice don't take 14/18 years to reach maturity like humans do);
They are also cheap, because they do not consume as humans do; There is also the simplicity factor, for example mice genomes are way simpler than humans ones,
and this helps on finding out how the genes work. And last, they have known geneome, I mean, the model organisms have sequenced their genomes and we know how they work.

I'm pretty sure that exists also other model organisms besides mice, like for bacteria (E.coli if I'm correct?).
We can see those model organism as a sort of heroes, because they help us understand how our body works, and how to treat diseases.

\vspace{0.5cm}
\subsection*{References}
BioInformatics' Lecture at FMI + what I know from my high school biology classes (2026). 